%!LW recipe=latexmk-xelatex
\documentclass[UTF8,zihao=-4]{ctexart}

% 页面与宏包
\usepackage[
  paperwidth=250mm, paperheight=245mm, margin=7mm]{geometry}
\usepackage{array}
\usepackage{fontspec}
\usepackage{multirow}

% 标题汉字：方正小标宋简体；正文汉字：仿宋_GB2312；数字字母：Times New Roman
\setmainfont{Times New Roman}
\setCJKmainfont[AutoFakeBold=2.2]{FangSong_GB2312}
\setCJKfamilyfont{xiaobiaosong}{FZXiaoBiaoSong-B05S}

% 页面与表格参数
\pagestyle{empty}
\setlength{\parindent}{0pt}
\setlength{\tabcolsep}{2.5pt}
\setlength{\arrayrulewidth}{0.45pt}
\renewcommand{\arraystretch}{1.18}
% 水平、垂直均居中的定宽列
\newcolumntype{C}[1]{>{\centering\arraybackslash}m{#1}}

% 文档信息与常用命令
\newcommand{\TableTitle}{北京航空航天大学 系号--学院/书院对照表}
\newcommand{\TableVersion}{2026.7}
\newcommand{\TitleFont}{\CJKfamily{xiaobiaosong}}

% 普通左右双栏行：系号、学院、系号、学院
\newcommand{\CollegePair}[4]{%
  #1 & \multicolumn{3}{c|}{#2} & #3 & \multicolumn{3}{c|}{#4} \\
  \hline
}

\begin{document}

\begin{center}
  % 标题严格居中；版本日期置于右上角且不参与标题居中计算
  \noindent
  \makebox[\linewidth][c]{{\zihao{2}\TitleFont\TableTitle}}%
  \makebox[0pt][r]{{\zihao{4}(\TableVersion)}}\\[2mm]

  % 八个基础列：普通学院行合并后三列；分组行直接使用基础列
  \begin{tabular}{|C{11mm}|C{13mm}|C{15mm}|C{68mm}|C{12mm}|C{15mm}|C{15mm}|C{70mm}|}
    \hline
    \bfseries 系号 & \multicolumn{3}{c|}{\bfseries 学院} &
    \bfseries 系号 & \multicolumn{3}{c|}{\bfseries 学院/书院} \\
    \hline

    % 普通学院
    \CollegePair{01}{材料科学与工程学院}{25}{国际学院}
    \CollegePair{02}{电子信息工程学院}{26}{新媒体艺术与设计学院}
    \CollegePair{03}{自动化科学与电气工程学院}{27}{化学学院}
    \CollegePair{04}{能源与动力工程学院}{28}{马克思主义学院}
    \CollegePair{05}{航空科学与工程学院}{29}{人文与社会科学高等研究院}
    \CollegePair{06}{计算机学院}{30}{空间与地球科学学院}
    \CollegePair{07}{机械工程及自动化学院}{31}{无人系统研究院}
    \CollegePair{08}{经济管理学院}{32}{航空发动机研究院}
    \CollegePair{09}{数学科学学院}{33}{体育部}
    \CollegePair{10}{生物与医学工程学院}{36}{国际前沿交叉科学研究院}

    % 37 系：北航学院及所属七个书院，仅在第 7～8 列之间绘制子行横线，保留合并单元格的完整性
    11 & \multicolumn{3}{c|}{人文社会科学学院（公共管理学院）} &
    \multirow{7}{*}{37} & \multirow{7}{*}{\shortstack{北\\航\\学\\院}} &
    71 & 传源书院（计算与智能科学类） \\
    \cline{1-4}\cline{7-8}

    12 & \multicolumn{3}{c|}{外国语学院} & & &
    73 & 士谔书院（信息科学与技术类） \\
    \cline{1-4}\cline{7-8}

    13 & \multicolumn{3}{c|}{交通科学与工程学院} & & &
    74 & 冯如书院（航空航天类） \\
    \cline{1-4}\cline{7-8}

    14 & \multicolumn{3}{c|}{可靠性与系统工程学院} & & &
    75 & 士嘉书院（航空航天类） \\
    \cline{1-4}\cline{7-8}

    15 & \multicolumn{3}{c|}{宇航学院} & & &
    76 & 守锷书院（航空航天类） \\
    \cline{1-4}\cline{7-8}

    16 & \multicolumn{3}{c|}{飞行学院} & & &
    77 & 致真书院（理科试验班类） \\
    \cline{1-4}\cline{7-8}

    17 & \multicolumn{3}{c|}{仪器科学与光电工程学院} & & &
    79 & 知行书院（社会科学试验班） \\
    \hline

    % 普通学院
    \CollegePair{18}{北京学院}{38}{医工交叉创新研究院}
    \CollegePair{19}{物理学院}{39}{网络空间安全学院}
    \CollegePair{20}{法学院}{41}{集成电路科学与工程学院}
    \CollegePair{21}{软件学院}{42}{人工智能学院（人工智能研究院）}
    \CollegePair{22}{继续教育学院}{45}{医学科学与工程学院}

    % 23 系：沈元学院及所属四个学院，横线从第 3 列开始，保留“23”和“沈元学院”的合并单元格
    \multirow{4}{*}{23} & \multirow{4}{*}{\shortstack{沈\\元\\学\\院}} &
    231 & 高等理工学院 & 54 & \multicolumn{3}{c|}{大科学装置研究院/零磁科学中心} \\
    \cline{3-8}

    & & 232 & 未来空天技术学院 & 57 & \multicolumn{3}{c|}{杭州国际创新研究院} \\
    \cline{3-8}

    & & 233 & 国家卓越工程师学院 & 60 & \multicolumn{3}{c|}{中法航空学院} \\
    \cline{3-8}

    & & 234 & 量子科技学院 & 61 & \multicolumn{3}{c|}{中法未来科技学院} \\
    \hline

    \CollegePair{24}{中法工程师学院/国际通用工程学院}{62}{中西智能学院}
  \end{tabular}
\end{center}

\end{document}
